\section{Introduction}
\label{sec:intro}

MQTT is the dominant application-layer protocol for constrained IoT
deployments. It is lightweight by design: a publish/subscribe broker, a
minimal binary header, and no mandatory authentication or transport security.
Those properties are what make it deployable on battery-powered sensors, and
they are also what make an MQTT network worth defending, because a broker that
accepts unauthenticated connections is a single point through which an attacker
can reach every subscriber. Network intrusion detection is consequently a
standard component of MQTT deployments, and the flow-level detector---one that
consumes summarised traffic records rather than packet payloads---is the form
that fits the bandwidth and privacy constraints of the setting.

Almost all such detectors treat a flow record as an independent feature
vector. A decision tree ensemble, a support vector machine or a multilayer
perceptron receives the packet counts, byte counts, inter-arrival statistics
and header flags of one flow and emits a label for that flow, with no
representation of the fact that the flow occurred between two particular
endpoints, or that those endpoints were simultaneously involved in hundreds of
other flows. The independence assumption is a modelling convenience, not a
property of the data.

Communication relationships plausibly carry information that the flow vector
does not. A port scan is not primarily a statement about any single flow; it is
a statement about one source contacting many destinations. A brute-force
attempt is a statement about repeated short exchanges with one service on one
host. Both patterns live in the structure of who talks to whom, and both are
only indirectly visible in the marginal statistics of an individual record.
Graph neural networks are the natural apparatus for such structure, and a
growing literature applies them to intrusion detection~\cite{bilot2024survey,
gnnids2026survey}, with E-GraphSAGE~\cite{lo2022egraphsage} establishing the
now-common recipe of treating flows as edges carrying attributes.

Two things about that literature motivate this paper. The first is that the
reported numbers leave no room to distinguish between methods. On
MQTT-IoT-IDS2020~\cite{hindy2020mqttdataset}, the dataset this paper uses, the
strongest published result is a genetic-algorithm-tuned bagging ensemble
reporting 99.997\% accuracy~\cite{norouzi2026reliable}. A method that reports
99.998\% has demonstrated nothing. When a benchmark is saturated, differences
in reported accuracy measure differences in preprocessing and split
construction rather than differences in detection capability, and the field has
repeatedly been warned about exactly this failure
mode~\cite{arp2024pitfalls,sommer2010outside}. Intrusion detection datasets
have also proven to carry labelling and construction artefacts substantial
enough to reverse published
conclusions~\cite{engelen2021troubleshooting,lanvin2022errors}.

The second is that a graph is not an unambiguous improvement over a feature
vector. Adding the endpoint pair to the model adds identity, and on a
laboratory testbed identity is a shortcut: the attacker occupies a fixed
address, so an edge incident to that address is an attack with near certainty.
A model can exploit this without learning anything about traffic behaviour, and
under a random split it will never be penalised for doing so. Shortcut learning
of this kind is well documented in machine learning
generally~\cite{geirhos2020shortcut}; what is missing for graph-based intrusion
detection is a protocol that detects it.

This paper therefore asks a narrower and more answerable question than whether
a graph helps. It asks what has to be true of an evaluation before the question
can be answered at all, and then answers it under those conditions. We
construct a communication multigraph in which IP endpoints are nodes and each
flow record is a distinct directed edge, so that every flow retains its own
attribute vector and receives its own prediction. We give nodes
behaviour-derived features computed from training traffic only, rather than
learned identity embeddings. We implement an edge-aware graph attention layer
that admits the flow attributes into the attention logit, and we compare it
against topology-only and edge-featured variants of the same architecture, a
reproduction of the base paper's detector, four classical learners, a
multilayer perceptron, and---critically---a baseline that sees nothing but the
endpoint pair.

Our contributions are the following. We characterise the evaluation
regime of MQTT-IoT-IDS2020 quantitatively, showing what a random split can and
cannot establish on it. We introduce endpoint-grouped evaluation protocols that
hold out whole devices, and we show which direction of holdout is well posed on
this testbed and which is not, with the structural reason. We contribute an
endpoint-identity baseline that requires no traffic features and that any
graph-based detector should be required to beat. We report a component
ablation that separates the contribution of graph topology, of flow attributes
in the message, and of flow attributes in the attention weight. We report
attention-based attribution, framed as attribution and not as explanation. And
we document, rather than discard, the conclusions that run against the
hypothesis we set out to test.

The remainder of the paper is organised as follows.
Section~\ref{sec:related} reviews graph-based and MQTT-specific intrusion
detection and situates the evaluation-methodology literature.
Section~\ref{sec:dataset} describes the dataset as released, including two
properties of its construction that determine how it must be labelled and
split. Section~\ref{sec:method} formalises the graph representation and the
edge-aware attention layer. Section~\ref{sec:methodology} sets out the
protocols, metrics and leakage controls. Section~\ref{sec:results} reports the
measurements, Section~\ref{sec:discussion} interprets them,
Section~\ref{sec:limitations} states what the study cannot support, and
Section~\ref{sec:conclusion} concludes.
